\section{Performance Evaluation}
\label{sec:performance}

\TW{} fixes the root/leaf transcript but leaves leaf scheduling to the
implementation. This section evaluates that separation: the same ciphertext and
tag are produced by scalar and vectorized implementations, while complete leaf
chunks can be batched to recover practical long-message throughput. The
comparison baseline is hardware-accelerated AES-128-GCM on the same hosts, so
the measurements should be read as an implementation sanity check and
same host reference point, not as a claim that \TW{} is faster than AES-GCM.

\subsection{Optimized Implementations}
\label{sec:impl}

We analyze a Go implementation of \TW{}\footnote{The implementation is available
at \url{https://github.com/codahale/treewrap}.} with optimized assembly paths
for the performance-critical permutation and absorb/squeeze loops. We
compare three vectorized assembly paths: \texttt{amd64} AVX-512,
\texttt{amd64} AVX2, and \texttt{arm64} NEON using the Armv8.2 SHA-3
extension. The framing, domain separation, padding, tree bookkeeping, and
dispatch remain in shared Go code. A pure Go path uses a scalar permutation
and a scalar eight-instance loop, and serves both as a reference path for
cross-checking the assembly cores and as a fallback on other architectures.

The tree structure of \TW{} exposes two distinct permutation workloads. The
root duplex---which absorbs the associated data, processes chunk~$0$, and
absorbs the leaf tag aggregation transcript---is inherently sequential, since
each call consumes the previous call's output. It is driven by a
\emph{single-duplex} ($\times 1$) core, except for its chunk-$0$ message phase.
The leaf duplexes operate on disjoint state and disjoint inputs
(\Cref{sec:tw128}), so complete leaf chunks are batched and run through an
\emph{eight-way} ($\times 8$) core. Leftover runs of two to seven complete
chunks use narrower remainder kernels, while a single leftover complete chunk,
a partial trailing chunk, and the count-aggregation tail fall back to the
single-duplex core.

Chunk~$0$ itself runs on the batched kernels rather than serially, a technique
we refer to as \emph{chunk~$0$ fusion}. The root duplex's chunk-$0$ message phase
has the same kernel-visible schedule as a leaf chunk and differs only in its
initial state. The implementation therefore seeds lane~$0$ of the first batched
call with the root state, runs chunk~$0$ alongside the first leaf chunks, then
writes lane~$0$'s output state back to the root duplex before the root absorbs
the leaf tags from that same call. This avoids a serial pass over chunk~$0$ and
is the reason the cycles-per-byte curve drops as soon as complete leaves are
present (\Cref{sec:throughput}). \Cref{app:kernels} gives the lane-major layout,
remainder-kernel strategies, lane~$0$ writeback details, and instruction-level
kernel structure.

\paragraph{Architecture-specific cores.}
On \texttt{amd64} the optimized implementation selects an AVX-512 or AVX2 core
once at startup from the \texttt{AVX512VL} CPUID feature bit. On \texttt{arm64}
the optimized implementation uses NEON with the Armv8.2 SHA-3 extension
(\texttt{FEAT\_SHA3}, available on Apple Silicon and Neoverse V1/V2 and N2).
Both architecture-specific paths provide a single-duplex core and batched leaf
kernels; the assembly stays below the mode boundary, so framing, domain
separation, padding, and tree bookkeeping remain shared with the reference path.

\paragraph{Data-independent execution.}
The scalar reference and all three vectorized implementations avoid
secret-dependent branches and secret-dependent memory or vector indices for the
key, nonce, associated-data contents, and plaintext or ciphertext contents. The
permutation and the XOR-based absorb/squeeze loops are straight-line with
respect to data values. The only indexed table is the round-constant array,
addressed by the public round number; dispatch, batched loads, stores, gathers,
scatters, and inactive-lane handling are controlled by public lengths, chunk
counts, and fixed chunk strides. \Cref{app:kernels} records the
architecture-specific addressing details.

\subsection{Benchmark Methodology}
\label{sec:bench-method}

We measure on two hosts, one per architecture. The \texttt{amd64} host is a
Google Compute Engine \texttt{c4-standard-4} instance with an Intel Xeon Platinum
8581C (Emerald Rapids, $2.30$\,GHz base, AVX-512 including \texttt{AVX512VL},
\texttt{AVX512\_VBMI2}, \texttt{GFNI}, and \texttt{VAES}), two physical cores
with simultaneous multithreading enabled, $48$\,KiB L1d and $32$\,KiB L1i per
core, $2$\,MiB L2 per core, and a $260$\,MiB shared L3, running Debian~12 (kernel
6.1) under KVM. The guest does not expose a \texttt{cpufreq} interface, so the
frequency governor and turbo state are not under our control; we mitigate the
resulting variance statistically (below). The \texttt{arm64} host is an Apple M4
Pro ($10$ performance cores and $4$ efficiency cores; per performance core,
$128$\,KiB L1i, $64$\,KiB L1d, and a $16$\,MiB shared L2) running macOS~26.5.1
(Darwin~25.5), with \texttt{FEAT\_SHA3} present. macOS exposes no userspace
frequency control.

Measurements are collected by a standalone harness (\texttt{cmd/cpb}) that locks
its goroutine to an OS thread and, for each algorithm, operation, and message
length, first calibrates an iteration count that fills at least $100$\,ms, then
records $100$ samples. It reports the median cycles per byte and median
wall-clock throughput together with the interquartile range as a dispersion
measure (\Cref{tab:dispersion}). Both metrics are read from the same timed
loop, so the cycles-per-byte and throughput figures for a given measurement
come from one interleaved run rather than two separately scheduled ones.

Each timed call uses a fixed all-zero key and nonce, empty associated data, and
preallocated source and destination buffers. Open measurements decrypt and
verify a ciphertext prepared before timing; seal and open reset the destination
slice length inside the timed loop but do not allocate per call. The reported
figures are therefore warm-cache, steady-state rates for message processing
under a fixed AEAD object. The single-threaded measurement leaves the SMT
sibling of the \texttt{amd64} host idle. All reported figures are
single-threaded: they exercise the SIMD axis of \TW{}'s parallelism, and
multicore scaling remains unmeasured. We report a sweep of seven message
lengths from $1$\,B to $16$\,MiB.

The cycle counters differ by architecture, and this affects how the
cycles-per-byte figures may be read. On \texttt{amd64} the harness reads
\texttt{RDTSC}; on this part the time-stamp counter is invariant and ticks at the
$2.30$\,GHz reference frequency, so the reported figures are \emph{reference}
cycles per byte and are unaffected by turbo (and convert directly to wall-clock
time). On \texttt{arm64} the harness reads \texttt{CNTVCT\_EL0} and scales it by
the ratio of the peak core frequency to \texttt{CNTFRQ\_EL0}; since Apple exposes
no frequency API, the peak is auto-detected by timing a dependent integer-add
chain and taking the top observed P-state, yielding \emph{core} cycles per byte.
Consequently the cycles-per-byte figures are comparable within an architecture
and against the same-host AES-128-GCM baseline, but not directly across
architectures; absolute throughput (\Cref{tab:throughput}) is the cross-platform
metric.

The baseline is Go's AES-128-GCM implementation
(\texttt{crypto/\allowbreak aes} with \texttt{crypto/\allowbreak cipher}).
It uses AES-NI and \texttt{PCLMULQDQ} on \texttt{amd64}, and the AES
and \texttt{PMULL} instructions on \texttt{arm64}.
The AES block cipher, GCM wrapper, and \TW{} AEAD object are constructed once
before timing, so the measurements exclude key schedule and GHASH key-power
precomputation. Each timed \TW{} call still performs the root duplex
initialization inside the call. Both directions are measured for each scheme;
the cycles-per-byte table reports both, and the throughput table reports the
encryption direction.

\subsection{Throughput}
\label{sec:throughput}

\Cref{tab:cpb} reports cycles per byte across the message-length sweep and
\Cref{tab:throughput} the corresponding wall-clock throughput in the bulk
regime. The main pattern in the \TW{} rows is the descent produced by the
tree-parallel leaf core. A message of at most one chunk has no complete leaf
chunk for chunk~$0$ to share a kernel with, so it runs entirely on the
single-duplex core, and the $8$\,KiB column sits at that core's rate. Past one
chunk, chunk~$0$ fusion (\Cref{sec:impl}) keeps every complete chunk on the
batched kernels, and the descent tracks lane occupancy.

On the \texttt{amd64} AVX-512 path, \TW{} encryption falls
$1.81 \to 0.73 \to 0.39$ cycles per byte across the $8$, $32$, and $64$\,KiB
columns and then holds at $0.38$--$0.39$ through $16$\,MiB. The AVX2 path has
the same shape at lower throughput: it runs the eight instances as two groups
of four rather than in a single $8$-wide register, leaving the AVX-512 core
about $2.9\times$ faster in bulk. The \texttt{arm64} curve flattens earlier
because the eight-way core processes its batch as four two-instance pairs; a
$32$\,KiB message, chunk~$0$ and three leaves as two full pairs, is already
close to the long-message rate. In wall-clock terms \TW{} reaches
$6.00$\,GB/s on the \texttt{amd64} AVX-512 host and $4.93$\,GB/s on the M4 Pro
for $16$\,MiB messages.

\begin{table}
\centering
\caption{Median cycles per byte ($100$ samples). The \texttt{amd64} figures are
\texttt{RDTSC} reference cycles at the $2.30$\,GHz invariant time-stamp counter;
the \texttt{arm64} figures are core cycles at the auto-detected peak performance-core
frequency. Cycle counts are not directly comparable across architectures
(\Cref{sec:bench-method}).}
\label{tab:cpb}
\begin{tabular}{l rrrrrrr}
\toprule
& \multicolumn{7}{c}{Message length} \\
\cmidrule(l){2-8}
& $1$\,B & $64$\,B & $8$\,KiB & $32$\,KiB & $64$\,KiB & $1$\,MiB & $16$\,MiB \\
\midrule
\multicolumn{8}{l}{\emph{Intel Xeon 8581C (Emerald Rapids), AVX-512}} \\
\quad \TW{} encrypt      & 571 & 9.04 & 1.81 & 0.73 & 0.39 & 0.38 & 0.38 \\
\quad \TW{} decrypt      & 605 & 9.55 & 1.82 & 0.73 & 0.39 & 0.38 & 0.38 \\
\quad AES-128-GCM seal   & 119 & 2.74 & 0.30 & 0.29 & 0.29 & 0.29 & 0.29 \\
\quad AES-128-GCM open   & 99 & 2.17 & 0.29 & 0.28 & 0.28 & 0.28 & 0.29 \\
\midrule
\multicolumn{8}{l}{\emph{Intel Xeon 8581C (Emerald Rapids), AVX2}} \\
\quad \TW{} encrypt      & 647 & 10.18 & 2.21 & 1.15 & 1.10 & 1.09 & 1.10 \\
\quad \TW{} decrypt      & 682 & 10.70 & 2.21 & 1.06 & 1.10 & 1.08 & 1.06 \\
\midrule
\multicolumn{8}{l}{\emph{Apple M4 Pro, NEON\,+\,\texttt{FEAT\_SHA3}}} \\
\quad \TW{} encrypt      & 620 & 9.88 & 1.97 & 0.92 & 0.89 & 0.89 & 0.91 \\
\quad \TW{} decrypt      & 671 & 10.65 & 2.04 & 0.90 & 0.87 & 0.87 & 0.89 \\
\quad AES-128-GCM seal   & 113 & 2.84 & 0.44 & 0.43 & 0.43 & 0.44 & 0.46 \\
\quad AES-128-GCM open   & 123 & 2.56 & 0.42 & 0.41 & 0.41 & 0.42 & 0.43 \\
\bottomrule
\end{tabular}
\end{table}

\begin{table}
\centering
\caption{Measured wall-clock throughput (GB/s, $\mathrm{GB} = 10^9$ bytes) from
the same \texttt{cmd/cpb} run as \Cref{tab:cpb}, encryption direction. The
\texttt{amd64} AES-128-GCM row is the AVX-512 host and is unaffected by the
\TW{} core selection.}
\label{tab:throughput}
\begin{tabular}{l rrrrr}
\toprule
& $8$\,KiB & $32$\,KiB & $64$\,KiB & $1$\,MiB & $16$\,MiB \\
\midrule
\multicolumn{6}{l}{\emph{Intel Xeon 8581C (Emerald Rapids)}} \\
\quad \TW{} (AVX-512)  & 1.27 & 3.17 & 5.93 & 6.10 & 6.00 \\
\quad \TW{} (AVX2)     & 1.04 & 2.00 & 2.10 & 2.10 & 2.08 \\
\quad AES-128-GCM      & 7.75 & 8.00 & 8.05 & 8.00 & 7.87 \\
\midrule
\multicolumn{6}{l}{\emph{Apple M4 Pro}} \\
\quad \TW{} (NEON+SHA3) & 2.27 & 4.84 & 5.04 & 5.01 & 4.93 \\
\quad AES-128-GCM       & 10.04 & 10.27 & 10.33 & 10.07 & 9.71 \\
\bottomrule
\end{tabular}
\end{table}

\paragraph{Dispersion.}
Each cell of \Cref{tab:cpb,tab:throughput} is the median of $100$ samples;
\Cref{tab:dispersion} summarizes their spread as the interquartile range
(IQR) relative to the median. The medians are stable: across the sweep the
relative IQR of the \TW{} cycles-per-byte measurements is at most $0.7\%$
on the \texttt{amd64} AVX-512 path, $0.9\%$ on AVX2, and $1.9\%$ on the M4
Pro, with per-path medians of $0.3\%$, $0.4\%$, and $1.2\%$. The full
min--max range is wider at isolated points---up to $4.9\%$, $12.0\%$, and
$16.7\%$ on the three paths---because turbo, governor, and scheduler states
are not under our control (\Cref{sec:bench-method}); these are tail
excursions that the IQR excludes. We therefore report medians throughout
and treat the IQR as the dispersion measure.

\begin{table}
\centering
\caption{Cycles-per-byte measurement dispersion for \TW{} across the
seven-length sweep (seal and open, $100$ samples per point), as the
interquartile range relative to the median. The final column is the widest
single-point min--max range; it is excluded from the IQR as a tail effect.}
\label{tab:dispersion}
\begin{tabular}{l rrr}
\toprule
& \multicolumn{2}{c}{Relative IQR} & Max \\
\cmidrule(lr){2-3}
Path & median & max & range \\
\midrule
Xeon 8581C, AVX-512 & $0.3\%$ & $0.7\%$ & $\phantom{0}4.9\%$ \\
Xeon 8581C, AVX2    & $0.4\%$ & $0.9\%$ & $12.0\%$ \\
Apple M4 Pro        & $1.2\%$ & $1.9\%$ & $16.7\%$ \\
\bottomrule
\end{tabular}
\end{table}

\subsection{Latency}
\label{sec:latency}

For short messages in this benchmark, the cost is set by the serial root
duplex and is independent of the parallel leaf machinery. With empty associated
data, a message of at most $\rho = 167$ bytes occupies a single rate block and
never enters leaf mode, so it costs exactly two permutations: one to initialize
the root duplex and one to close the message block and extract the tag. The
$1$\,B and $64$\,B columns of \Cref{tab:cpb} bear this out: per-message latency
is nearly flat across them at roughly $571$ reference cycles on \texttt{amd64}
($571$ at $1$\,B versus $9.04 \times 64 \approx 578$ at $64$\,B) and $620$
core cycles on \texttt{arm64} ($620$ versus $9.88 \times 64 \approx 633$),
i.e.\ about $285$ and $310$ cycles per \Keccakp{} respectively. Because such
messages bypass the leaf core entirely, leaf-level parallelism imposes no
latency penalty on them.

The latency/throughput trade-off appears instead in the intermediate regime,
which chunk~$0$ fusion has narrowed to roughly the one-to-eight-chunk band. The
worst per-byte cost in the sweep is the single-chunk regime at $8$\,KiB
($1.8$--$2.2$ cycles per byte across the three paths), where the message is the
chunk-$0$ phase alone on the single-duplex core. Once the message has complete
leaf chunks for chunk~$0$ to share a kernel with, the cost falls quickly: by
$32$\,KiB \TW{} is within a factor of two of its AVX-512 floor and close to its
\texttt{arm64} floor, and the floor itself is reached once the eight-way core
fills at $64$\,KiB. The root duplex also bounds the tail of a long
message: the final tag cannot be produced until chunk~$0$ has been processed and
every leaf tag has been absorbed into the aggregation transcript. That absorption
is cheap---each rate block of the root absorbs just over five $32$-byte leaf tags, so a
$16$\,MiB message ($\approx 2000$ leaves) adds on the order of $400$ root
permutations against roughly $10^5$ permutations of leaf work, under $0.5\%$
overhead---so the serial root does not materially cap bulk throughput.

\subsection{AES-128-GCM Baseline}
\label{sec:compare-perf}

We compare against the one baseline measured under the same harness on the same
hosts, the Go standard library's hardware-accelerated AES-128-GCM
(\Cref{tab:cpb,tab:throughput}). The harness times \TW{} and AES-128-GCM
back-to-back at each length and reads both metrics from the same loop, so the
cycles-per-byte and throughput comparisons hold under identical conditions and
agree. In the bulk regime AES-128-GCM is faster on both architectures, as
expected for a primitive with dedicated AES and carryless-multiply instructions:
at $16$\,MiB it runs at $0.29$ cycles per byte on \texttt{amd64} and $0.46$ on
\texttt{arm64}, against $0.38$ and $0.91$ for \TW{}---about $1.3\times$ and
$2.0\times$ faster. The wall-clock throughputs show the same margins: \TW{}
reaches $6.00$ and $4.93$\,GB/s against AES-128-GCM's $7.87$ and
$9.71$\,GB/s. The wider \texttt{arm64} gap reflects the strength of Apple's AES
and \texttt{PMULL} units. \TW{} attains these rates with a permutation that has
no dedicated hardware support, relying only on the SHA-3 extension on
\texttt{arm64} and on general SIMD on \texttt{amd64}.

AES-128-GCM is also faster for short messages in this steady-state measurement:
at $1$\,B it costs $119$ and $113$ cycles on the two hosts against \TW{}'s
$571$ and $620$. This is the expected result when AES-GCM key setup is
amortized outside the timed call. In the sub-bulk regime AES-128-GCM is several
times faster at one chunk ($0.30$ versus $1.81$ cycles per byte at $8$\,KiB on
\texttt{amd64}), but the gap narrows once the fused kernels engage: about
$2.5\times$ at $32$\,KiB and the bulk ratio from $64$\,KiB on. \TW{}'s intended
value relative to AES-128-GCM lies in the tag-centric hash and committing
properties analyzed in \Cref{sec:root-tag-hash-security}, with the baseline
AEAD bounds of \Cref{sec:aead-security}, rather than in raw speed.
